\documentclass[sigconf]{acmart}

% uncomment to enable double-spacing
%\usepackage[doublespacing]{setspace}
\usepackage{booktabs} 
\usepackage{hyperref}
\usepackage[numbib]{tocbibind}
\renewcommand\bibname{References}

% Copyright
\setcopyright{rightsretained}
\copyrightyear{2018}

% Formatting
\settopmatter{printacmref=false} 
\renewcommand\footnotetextcopyrightpermission[1]{} 
\pagestyle{plain} 

\begin{document}
\title{TDA Mapper Paper Template}

\author{John Q Public}
\affiliation{
  \institution{University of Iowa}
  \streetaddress{14 MacLean Hall}
  \city{Iowa City}
  \state{Iowa}
  \postcode{52242-1419}
}
\email{john-q-public@uiowa.edu}

% Use this if the default list of authors is too long for headers.
%\renewcommand{\shortauthors}{B. Trovato et al.}

\maketitle

\section{Abstract}
TODO

\section{Introduction}
%Briefly introduce the problem, techniques, and outline the paper. Try to use as few technical terms as possible (or reference section where defined).

TODO

\section{Background}
%Problem description: Fully describe the problem. Describe how the data is created, what is its format, what are issues that one should consider (for example are their different types of noise), etc. Please keep in mind that people from a variety of backgrounds may be interested in your article, so please help them understand your data. Deeper mathematics can be applied if more people have a better understanding of the problem.

%Mathematical background: One can reference appropriate papers that describe the methods used in your paper or you can provide the background yourself. Including motivation specific to your problem would be particularly helpful.

TODO

\begin{figure}[h]
\centering
\includegraphics[width=\columnwidth]{sample-figure.png}
\caption{Sample figure}
\label{sample-figure}
\end{figure}

Using Figure \ref{sample-figure}, we see...

\section{Results}
%Data analysis: In either the results or discussion section, motivate your choice of software. Software often requires one to choose various values for different parameters. Motivate your parameter choices. Are these choices robust (for example, does one get similar results for different choices of parameter values).

TODO

\section{Discussion}
%What do your results mean?: This will depend on your application. For example, if you find a hole in your data, you might be able to determine exactly where that hole comes from as was done in Adam's and Carlsson's On the nonlinear statistics of range image patches.

%Compare your results to results using other types of data analysis tools.

%Validation: Part of validation includes motivating and determining the robustness of the choices you made when analyzing the data and comparing your results to previously published results. But another part is to determine if the results give you useful information for your particular problem. For example, a cluster found in Topology based data analysis identifies a subgroup of breast cancers with a unique mutational profile and excellent survival by Nicolau, Levine, and Carlsson was validated by further examining the biological properties of this cluster. Of course, this requires laboratory experiments which your collaborators may or may not be able to perform.

TODO

\section{Conclusions}
TODO

\section{Acknowledgements}
TODO

\section{Author Contributions}
TODO

\section{Funding Sources and Conflicts of Interest}
TODO

\bibliographystyle{ACM-Reference-Format}
\bibliography{bibliography}

\end{document}